\documentclass[11pt]{article}
\usepackage[margin=1in]{geometry}
\usepackage{amsmath,amssymb,amsthm,xcolor,hyperref}
\hypersetup{colorlinks=true,allcolors=blue!45!black}
\newtheorem{theorem}{Theorem}
\newtheorem{corollary}[theorem]{Corollary}
\title{A Gram Certificate for Monochromatic-Box Matrices}
\author{Myeongjin Kang}
\date{October 25, 2026}
\begin{document}
\maketitle
\begin{abstract}
We formalize positive semidefiniteness of the monochromatic-box matrix by an exact Gram identity. The proof applies to arbitrary finite color sets and arbitrary subset sizes. The original conjecture was already resolved by Cooper and Forst in 2017, including the general set-intersection Gram argument. This is a verification note, not a claim of a newly solved open problem.
\end{abstract}
\section{The matrix and its status}
For $r\ge0$, let $\mathcal F$ be the set of binary colorings of an $r$-element set. The source matrix is
\[
M_r(f,g)=\binom{|f^{-1}(0)\cap g^{-1}(0)|}{2}
 +\binom{|f^{-1}(1)\cap g^{-1}(1)|}{2}.
\]
The change from colors $1,2$ in \cite{CFP} to colors $0,1$ only relabels indices. This is Definition 7.2 in the original preprint, followed by Conjecture 7.4. In the published version the corresponding labels are Definition 16 and Conjecture 18. Cooper and Forst's Corollary 3 \cite{CF} resolves the positivity question for every $r\ge0$; Proposition 4 supplies the underlying intersection-matrix argument. The MathDB record retrieved for this project still calls it open, and is therefore outdated.
\section{A general incidence identity}
\begin{theorem}\label{thm:gram}
Let $A,C$ be finite sets and let $k\ge0$. Index a real matrix $M$ by functions $f:A\to C$, with entries
\[
M(f,g)=\sum_{c\in C}\binom{|\{a\in A:f(a)=g(a)=c\}|}{k}.
\]
Then $M$ is positive semidefinite.
\end{theorem}
\begin{proof}
For a coloring $f$, a color $c$, and a subset $T\subseteq A$, put
\[
B(f,(c,T))=
\begin{cases}
1,& |T|=k\text{ and }f(a)=c\text{ for all }a\in T,\\
0,&\text{otherwise}.
\end{cases}
\]
This is a finite real matrix. Fix $f,g,c$. The product
$B(f,(c,T))B(g,(c,T))$ equals one precisely when $T$ is a $k$-element subset of the common fiber $\{a:f(a)=g(a)=c\}$, and is zero otherwise. Thus summing over all subsets $T$ gives the binomial coefficient in the definition of $M$. Summing over colors yields $M=BB^{\mathsf T}$. It follows that, for every real vector $v$ indexed by the colorings,
\[
v^{\mathsf T}Mv=\|B^{\mathsf T}v\|^2
 =\sum_{c,T}\left(\sum_fB(f,(c,T))v_f\right)^2\ge0.
\]
The factorization also proves that $M$ is symmetric, completing positive semidefiniteness.
\end{proof}
\begin{corollary}
The source matrix $M_r$ is positive semidefinite for every nonnegative integer $r$, and in particular for all $r\ge3$ as required by the original conjecture.
\end{corollary}
\begin{proof}
Apply Theorem~\ref{thm:gram} with $A=\{0,\ldots,r-1\}$, $C=\{0,1\}$, and $k=2$.
\end{proof}
The proof remains valid when $k>|A|$, when all entries vanish, or when $k=0$, when the empty subset contributes once per color. No exceptional size assumption is needed.
\section{Formal coverage and reproducibility}
The file \path{formal/Monochromatic.lean} defines the common fibers, the binomial-entry matrix, and the incidence matrix. The declaration \path{feature_count} reduces incidence counts to \path{Finset.card_powersetCard}; \path{gram_identity} proves the factorization. Positive semidefiniteness then follows from the standard library theorem for a matrix times its conjugate transpose. Over the reals this transpose is the usual transpose. The declarations \path{sourceMatrix} and \path{sourceMatrix_eq_boxMatrix} explicitly recover the source's two binomial terms, and \path{mathdb317856} is the exact specialized conclusion.

The rebuild driver uses Lean 4.34.1 and mathlib revision \path{d13f23b7} (full hash in the supplied pins). It checks every definition and theorem, removes stale local objects, treats warnings as errors, and enforces an axiom allowlist containing only propositional extensionality, classical choice, and quotient soundness. The packet contains versioned source snapshots and SHA-256 manifests. It does not formalize or assert the separate spectral formulas in the cited papers.
\begin{thebibliography}{9}
\bibitem{CFP} J. Cooper, S. Fenner, and S. Purewal, \emph{Monochromatic Boxes in Colored Grids}, arXiv:0810.3019v1 (2008), Definition 7.2 and Conjecture 7.4. \url{https://arxiv.org/html/0810.3019v1}.
\bibitem{CF} J. Cooper and M. Forst, \emph{Positive Semidefiniteness of Matrices arising from Ramsey Theory}, arXiv:1704.01644v2 (2017), Corollary 3 and Proposition 4. \url{https://arxiv.org/html/1704.01644v2}.
\end{thebibliography}
\end{document}
